\documentclass[9pt,twocolumn]{pnas-new}

\templatetype{pnasresearcharticle}

\title{Metabolic niche partitioning drives tri-trophic coexistence in cross-feeding microbial communities}

\author[a]{Jian Wang}

\affil[a]{Department of Mathematics and Biological Sciences, [Institution Name]}

\leadauthor{Wang}

\significancestatement{Microbial communities exhibit remarkable metabolic diversity, yet the mechanisms enabling coexistence of specialists and generalists remain poorly understood. By combining analytical theory with systematic parameter space analysis, we demonstrate that metabolic flexibility—the ability to switch between resource utilization pathways—creates ecological niches that stabilize complex communities. Our mathematical framework reveals precise thresholds governing community assembly and identifies transcritical bifurcations as a universal mechanism for generalist invasion. These findings provide quantitative guidelines for engineering synthetic consortia and predict testable phase transitions in natural communities responding to environmental change. The parameter landscapes we characterize offer a roadmap for understanding and manipulating the metabolic architecture of microbial ecosystems.}

\authorcontributions{J.W. designed research, performed research, analyzed data, and wrote the paper.}

\authordeclaration{The author declares no competing interest.}

\correspondingauthor{\textsuperscript{1}To whom correspondence should be addressed. E-mail: jian.wang@institution.edu}

\begin{abstract}
Cross-feeding interactions, where metabolic by-products of one species become resources for another, structure microbial communities from the human gut to ocean biofilms. Understanding how metabolically flexible generalists coexist with specialized cross-feeders has profound implications for community stability, function, and engineering. Here, we develop a mathematical framework analyzing a three-species system comprising substrate specialists, metabolite specialists, and pathway-switching generalists. Through complete analytical characterization and systematic parameter space mapping, we identify a transcritical bifurcation governing generalist invasion: community composition transitions sharply at critical pathway utilization thresholds. The invasion boundary in parameter space reveals that generalists must occupy intermediate metabolic strategies—neither fully substrate- nor metabolite-specialized—to coexist with specialists. This niche partitioning through metabolic flexibility stabilizes species diversity and creates testable predictions for experimental systems. Our parameter landscape analysis demonstrates how pairwise mutualistic strength modulates higher-order assembly rules, offering design principles for synthetic biology and conservation strategies for natural communities facing metabolic perturbations.
\end{abstract}

\dates{This manuscript was compiled on \today}

\begin{document}

\maketitle
\thispagestyle{firststyle}

\section*{Introduction}

Microorganisms dominate Earth's biogeochemical cycles, yet despite intense competition for resources, microbial ecosystems consistently harbor tremendous diversity \cite{hug2016new,louca2016decoupling}. This paradox—high diversity under apparent competitive exclusion—has driven decades of ecological theory seeking coexistence mechanisms \cite{hutchinson1961paradox,hardin1960competitive}. Cross-feeding interactions, wherein metabolic waste products of one species serve as essential nutrients for another, represent a fundamental pathway toward stable coexistence \cite{ponomarova2015metabolic,goldford2018emergent}. Such syntrophic partnerships underpin critical processes from wastewater treatment to human health \cite{embree2015networks,harcombe2014metabolic}, yet theoretical understanding of community assembly rules in cross-feeding systems remains incomplete.

Experimental evolution studies have demonstrated that metabolic cross-feeding readily evolves in both eukaryotic \cite{shou2007synthetic,waite2015adaptation} and prokaryotic \cite{rosenzweig2013microbial,mccully2017emergence} systems, often accompanied by diversification into specialist and generalist strategies \cite{garcia2021eco,goldford2018emergent}. Specialists, optimized for narrow resource niches, coexist with generalists capable of utilizing multiple pathways, creating metabolic division of labor \cite{harcombe2014metabolic,ziesack2019engineered}. However, the precise conditions enabling such coexistence—particularly when generalist metabolic strategies can be tuned continuously—lack rigorous mathematical characterization. This gap is critical because engineering stable synthetic consortia requires predicting assembly outcomes across parameter space \cite{kong2018designing,liu2018coupling}, while understanding natural community responses to environmental perturbations demands knowledge of bifurcation thresholds separating alternative stable states \cite{scheffer2001catastrophic,dai2013niche}.

Here we develop a complete analytical framework for a paradigmatic three-species cross-feeding system. Our model captures substrate specialists that produce metabolites, obligate metabolite specialists requiring cross-fed nutrients, and metabolically flexible generalists that can partition effort between substrate and metabolite utilization pathways. By introducing net interaction parameters that linearly depend on a pathway-switching parameter $\omega$, we achieve full analytical tractability while retaining biological realism. This formulation enables systematic parameter space mapping that reveals how mutualistic strength, competitive intensity, and metabolic strategy jointly determine community structure.

Our analysis yields three major advances. First, we derive explicit formulas for invasion thresholds, demonstrating that generalist establishment requires pathway utilization falling within a bounded window—too substrate-specialized and generalists compete with substrate specialists; too metabolite-specialized and they displace obligate mutualists. Second, we identify transcritical bifurcations as the assembly mechanism, wherein equilibrium densities transition smoothly through critical points where community membership changes. Third, through comprehensive parameter landscape visualization, we demonstrate how pairwise interaction strengths propagate to constrain three-species assembly, providing quantitative design principles for community engineering. These findings bridge theoretical ecology and synthetic biology, offering both fundamental insights into niche partitioning mechanisms and practical guidelines for manipulating microbial ecosystems.

\section*{Results and Discussion}

\subsection*{Obligate mutualism establishes a stable two-species platform}

\textbf{Key Finding:} \textit{Obligate cross-feeding creates quantitative thresholds that constrain synthetic syntrophy and predict natural methanogenic partnerships.}

We begin by analyzing the foundational two-species subsystem comprising substrate specialists (S) and metabolite specialists (M). The scaled dynamics are:
\begin{align}
\frac{ds}{dt} &= r_S s(1 + \sigma_{SM} m - s) \label{eq:SM1}\\
\frac{dm}{dt} &= r_M m(-1 + \sigma_{MS} s - m) \label{eq:SM2}
\end{align}
where $s$ and $m$ denote population densities scaled by carrying capacities, $r_S$ and $r_M$ are intrinsic growth rates, $\sigma_{SM}$ quantifies benefit to S from M's metabolic by-products, and $\sigma_{MS}$ measures M's dependence on S-provided substrates. Critically, M exhibits a negative basal growth rate (the $-1$ term in Eq.~\ref{eq:SM2}), representing obligate cross-feeding: M cannot survive independently and requires $\sigma_{MS} s > 1$ for positive growth.

Analytical solution of the equilibrium conditions reveals a unique coexistence state at $s^* = (1-\sigma_{SM})/(1-\sigma_{MS}\sigma_{SM})$ and $m^* = (\sigma_{MS}-1)/(1-\sigma_{MS}\sigma_{SM})$, existing when $\sigma_{MS} > 1$ (M viability) and $\sigma_{MS}\sigma_{SM} < 1$ (stability constraint). Jacobian analysis confirms stability through negative trace and positive determinant, yielding the combined criterion: mutualism must be sufficiently strong for M survival but not so strong as to create runaway positive feedback. This S-M consortium serves as the ecological platform upon which generalist invasion occurs.

The biological interpretation is profound. Natural obligate syntrophies—from methanogens requiring hydrogen-producing fermenters \cite{schink2002syntrophic,stams2006electron} to auxotrophic bacterial pairs exchanging amino acids \cite{d2008auxotroph,pande2015fitness}—must satisfy these quantitative thresholds. The parameter $\sigma_{MS} > 1$ represents a minimum cross-feeding flux, experimentally tunable through nutrient availability or genetic manipulation of secretion pathways \cite{mee2014syntrophic,harcombe2014metabolic}. The stability bound $\sigma_{MS}\sigma_{SM} < 1$ predicts an upper limit to reciprocal benefit, beyond which oscillations or extinction occur. Empirical measurements in engineered yeast pairs \cite{shou2007synthetic} and bacterial consortia \cite{hoek2016resource} support such boundaries, though systematic mapping across parameter space remains experimentally challenging. Our analytical solution provides testable predictions: varying inducer concentrations controlling cross-feeding gene expression should map a phase diagram matching our stability criteria.

\subsection*{Generalist invasion requires intermediate metabolic strategies}

\textbf{Key Finding:} \textit{Metabolic flexibility becomes advantageous only at intermediate pathway allocation—extreme specialists exclude generalists, resolving the paradox of specialist-generalist coexistence in cross-feeding communities.}

The generalist (G) can utilize both substrate and metabolite pathways, with relative investment controlled by parameter $\omega \in [0,1]$. At $\omega = 0$, G behaves identically to M (pure metabolite specialist); at $\omega = 1$, G resembles S (pure substrate specialist); intermediate values represent true generalism. This generates net interactions:
\begin{align}
a &= (1-\omega)\sigma_{SG} - \omega\alpha_{SG} \quad \text{(G effect on S)}\\
c &= (1-\omega)\sigma_{GS} - \omega\alpha_{GS} \quad \text{(S effect on G)}\\
d &= 2\omega - 1 \quad \text{(G basal growth rate)}
\end{align}
where $\sigma$ terms denote cooperative benefits and $\alpha$ terms competitive costs. Crucially, all parameters are linear in $\omega$, enabling analytical tractability.

To determine when G invades the stable S-M equilibrium, we compute G's invasion fitness at $(s^*_{SM}, m^*_{SM}, 0)$. The per-capita growth rate of rare generalists is $\lambda_G = r_G(d + c s^*_{SM} + e m^*_{SM})$, where $e = \omega\sigma_{GM} - (1-\omega)\alpha_{GM}$ represents G's interaction with M. Invasion succeeds when $\lambda_G > 0$, yielding a critical threshold:
\begin{equation}
\omega_{crit} = \frac{1 - \sigma_{GS} s^*_{SM} + \alpha_{GM} m^*_{SM}}{2 - (\sigma_{GS}+\alpha_{GS}) s^*_{SM} + (\sigma_{GM}+\alpha_{GM}) m^*_{SM}} \label{eq:omegacrit}
\end{equation}

This explicit formula—our first major result—reveals how pairwise S-M structure propagates to constrain generalist assembly. The numerator contains baseline pathway viability ($2\omega - 1$) modulated by S-mediated facilitation and M-mediated competition. The denominator combines interaction strengths, effectively weighting pathway switching costs. For our baseline parameters ($\sigma_{MS} = 1.5$, $\sigma_{SM} = 0.5$, $\sigma_{GS} = 0.4$, $\alpha_{GM} = 0.3$), Eq.~\ref{eq:omegacrit} predicts $\omega_{crit} \approx 0.35$. Below this threshold, generalists are too metabolite-specialized, competing directly with M for cross-fed resources and facing exclusion. Above threshold, substrate pathway investment creates sufficient niche differentiation for coexistence.

The ecological significance extends beyond mathematical elegance. Metabolic flexibility—often considered advantageous—here becomes a double-edged sword: only intermediate strategies enable coexistence. This resolves the apparent paradox of specialists and generalists coexisting in cross-feeding systems \cite{rosenfeld2020defining,goldford2018emergent}. Generalists avoid direct competition not through superior efficiency but through niche occupancy unavailable to specialists. Experimental tests could employ inducible promoters controlling pathway gene expression \cite{guantes2006multistable}, varying inducer concentrations to scan $\omega$ and measuring community composition. The predicted sharp transition at $\omega_{crit}$ should manifest as sudden generalist establishment, analogous to range expansion thresholds in spatial ecology \cite{hallatschek2007genetic}.

\subsection*{Parameter space architecture governs assembly outcomes}

\textbf{Key Finding:} \textit{Curved invasion boundaries reveal nonlinear coupling between mutualistic strength and metabolic strategy, creating predictable fragile zones where communities are sensitive to nutrient perturbations versus robust regions buffered against environmental change.}

To understand community assembly more broadly, we systematically mapped invasion boundaries across two-dimensional parameter spaces (Fig.~2). Panel A displays the invasion rate $\lambda_G$ across the $(\omega, \sigma_{MS})$ plane. The zero-contour (thick black line) partitions parameter space into exclusion (red, $\lambda_G < 0$) and invasion (green, $\lambda_G > 0$) regions. Strikingly, the boundary is neither linear nor parallel to axes, instead exhibiting curvature that reflects nonlinear coupling between pathway strategy and mutualistic strength.

Several critical features emerge. First, at low $\sigma_{MS}$ (weak S-to-M benefit), generalist invasion requires higher $\omega$—stronger substrate specialization compensates for weak mutualistic support. Conversely, strong mutualism ($\sigma_{MS} > 2$) permits generalist establishment even at low $\omega$, suggesting highly mutualistic platforms readily accommodate diverse metabolic strategies. Second, the invasion boundary asymptotically approaches $\omega \approx 0.3$ as $\sigma_{MS} \to \infty$, indicating an intrinsic minimum substrate utilization necessary for generalist viability regardless of mutualistic context. Third, the gradient magnitude varies across parameter space: shallow gradients near the boundary imply high sensitivity to perturbations, whereas steep regions confer robustness.

These patterns have profound implications for natural and engineered systems. Communities experiencing nutrient shifts that alter cross-feeding fluxes (modeled as changing $\sigma_{MS}$) will undergo predictable composition changes. For instance, reducing carbon substrate availability in soil communities should favor generalists with higher $\omega$ (more substrate-specialized), while environments rich in cross-fed metabolites (high $\sigma_{MS}$) tolerate broader metabolic diversity. Synthetic consortium design can exploit these landscapes: to ensure generalist inclusion, engineer host-producer pairs with $\sigma_{MS}$ exceeding the threshold for target $\omega$, or tune inducible pathway expression to match the predicted invasion window.

Panels B and C reveal mechanistic underpinnings by showing S-M equilibrium densities across parameter space. The S-M equilibrium provides the "ecological landscape" into which generalists invade. Higher $\sigma_{MS}$ elevates $m^*_{SM}$ (panel C), creating more metabolite-mediated facilitation opportunities for generalists (captured in the $e m^*_{SM}$ term of $\lambda_G$). Simultaneously, $s^*_{SM}$ decreases with $\sigma_{MS}$ (panel B), reducing substrate competition. These opposing trends create the curved invasion boundary: generalist success requires balancing substrate competition against metabolite facilitation, with the optimal $\omega$ shifting as mutualistic structure changes.

Panel D examines the $(\omega, \sigma_{SM})$ plane holding $\sigma_{MS}$ fixed, revealing a complementary constraint. The horizontal dashed line marks $\sigma_{SM} = 1/\sigma_{MS}$, the S-M stability threshold. Above this line, S-M mutualism becomes unstable (runaway growth), collapsing the platform for generalist invasion. This upper bound demonstrates a hierarchy in assembly constraints: pairwise stability must precede higher-order invasion. Even strongly substrate-specialized generalists ($\omega \to 1$) cannot invade if S-M destabilizes first. This finding informs conservation: preserving syntrophic partnerships requires maintaining parameter regimes supporting both pairwise stability and invasion windows for additional diversity.

Panel E explores cooperation-competition balance in S-G interactions. The zero-invasion contour curves through $(\sigma_{GS}, \alpha_{GS})$ space, crossing the diagonal where cooperation equals competition. Generalist invasion requires cooperative benefits exceeding competitive costs, but the threshold is not simply $\sigma_{GS} > \alpha_{GS}$. Instead, M-mediated effects (the $e m^*_{SM}$ term) shift the boundary, illustrating how three-species interactions are irreducible to pairwise sums. This invalidates simplistic "net interaction" heuristics: predicting invasion demands accounting for all species' influences mediated through equilibrium community structure.

\subsection*{Transcritical bifurcations mediate community assembly}

\textbf{Key Finding:} \textit{Sharp community transitions at critical thresholds—structurally stable transcritical bifurcations—enable quantitative prediction of composition changes in response to metabolic perturbations, analogous to ecosystem tipping points but mathematically tractable.}

Systematically varying $\omega$ from 0 to 1 traces a continuous trajectory through parameter space, revealing the bifurcation structure underlying community transitions. Below $\omega_{crit}$, the S-M equilibrium is stable and attracts all trajectories; generalists, if introduced, decline to extinction. At $\omega = \omega_{crit}$, a transcritical bifurcation occurs: the S-M equilibrium loses stability along the generalist direction, and an interior equilibrium (three-species coexistence) emerges from the boundary. For $\omega > \omega_{crit}$, the interior equilibrium is stable, and generalists persist. This smooth exchange of stability—characteristic of transcritical bifurcations—contrasts with saddle-node bifurcations (where equilibria appear or disappear discontinuously) or Hopf bifurcations (generating oscillations).

The transcritical mechanism has deep ecological significance. Community assembly is often conceptualized as sequential invasion events \cite{law1996permanence,urban2008landscape}, but the dynamical underpinnings remain underexplored. Our analysis demonstrates that invasion transitions can exhibit mathematical signatures (zero eigenvalues at bifurcation points, smooth equilibrium branches) enabling quantitative prediction. Moreover, transcritical bifurcations are structurally stable: small parameter perturbations shift $\omega_{crit}$ but preserve the bifurcation type. This robustness suggests that experimental systems—prone to environmental fluctuations—should still exhibit the predicted invasion threshold, albeit with noise-induced variability near the boundary.

A second bifurcation occurs at higher $\omega$ values (denoted $\omega_{crit2}$), where the metabolite specialist M is displaced. As $\omega$ increases, generalists increasingly specialize on substrates, competing with S while reducing metabolite dependence. Eventually, generalist-mediated changes to resource availability make M unviable. The three-species equilibrium merges with an S-G boundary equilibrium, and M density smoothly declines to zero—another transcritical bifurcation. Thus, the generalist-invaded community exists only for $\omega \in (\omega_{crit1}, \omega_{crit2})$, defining a coexistence window. Outside this interval, the community is two-species (S-M at low $\omega$, S-G at high $\omega$), highlighting how continuous metabolic tuning can trigger discrete composition changes.

These findings address a longstanding question in microbial ecology: how do communities maintain functional redundancy (multiple species performing similar roles) alongside functional specialization? Our model suggests a resolution. Within the coexistence window, specialists and generalists partition metabolic space, with generalists occupying an intermediate niche inaccessible to specialists constrained to pathway extremes. This metabolic niche partitioning complements traditional resource partitioning, potentially explaining microbial diversity in environments where classical resource competition theory predicts exclusion \cite{wilson1999there,scheffer2018towards}.

\subsection*{Implications for engineered and natural systems}

The analytical framework and parameter landscapes developed here offer actionable guidance for synthetic consortium design. First, choose specialist pairs satisfying $\sigma_{MS} > 1$ and $\sigma_{MS}\sigma_{SM} < 1$, ensuring a stable mutualistic platform. This can be achieved genetically by engineering cross-feeding pathways with defined secretion rates \cite{mee2014syntrophic,harcombe2014metabolic} or ecologically by selecting species with complementary metabolic capabilities \cite{ziesack2019engineered}. Second, to incorporate generalists, employ inducible promoters controlling pathway allocation and tune inducer concentrations to place $\omega$ within the coexistence window. Our explicit formulas (Eq.~\ref{eq:omegacrit} and parameter maps Fig.~2) enable a priori prediction of permissive inducer ranges, reducing trial-and-error experimentation.

For natural systems, our results predict testable responses to environmental perturbations. Nutrient amendments that increase cross-feeding fluxes (raising $\sigma_{MS}$) should expand the coexistence window, admitting generalists with broader metabolic strategies. Conversely, resource limitation reducing cross-feeding should narrow the window, excluding generalists unless they shift $\omega$ toward extreme specialization. Time-series monitoring of community composition during nutrient transitions could empirically validate predicted bifurcation thresholds. Additionally, the parameter sensitivity revealed in gradient maps (Fig.~2) identifies fragile versus robust regions: communities near the invasion boundary exhibit high compositional variability under environmental noise, whereas those deep within invasion or exclusion zones resist perturbations. This could inform ecosystem management strategies, prioritizing conservation efforts on communities in sensitive parameter regimes.

The broader implications extend to evolutionary dynamics. If $\omega$ itself evolves—through mutations in pathway regulation—then adaptive dynamics theory \cite{geritz1998evolutionarily} predicts evolutionary branching when disruptive selection favors extreme strategies. Our bifurcation structure suggests this occurs if mutualistic strength falls outside stability bounds, driving generalists toward specialist extremes. Conversely, within the coexistence window, stabilizing selection maintains intermediate $\omega$, preserving generalism. Experimental evolution studies allowing pathway regulation to mutate could test whether predicted evolutionary endpoints match analytical bifurcation diagrams \cite{good2017dynamics,lee2019rate}.

Finally, our approach demonstrates the power of simplified, analytically tractable models for generating mechanistic insight. While natural communities involve dozens to thousands of species, the three-species system captures essential features—obligate mutualisms, metabolic flexibility, niche partitioning—that likely recur in more complex assemblages. The parameter reduction strategy, collapsing cooperation-competition trade-offs into net interactions linear in a control parameter, could extend to larger systems, potentially revealing universal assembly rules governing microbial metacommunities.

\section*{Conclusion}

By integrating complete analytical characterization with systematic parameter space exploration, we have demonstrated how metabolic niche partitioning stabilizes tri-trophic coexistence in cross-feeding communities. The identification of transcritical bifurcations as assembly mechanisms, derivation of explicit invasion thresholds, and mapping of multidimensional parameter landscapes provide both fundamental insights and practical tools. Our work establishes quantitative links between molecular-level pathway regulation (encoded in $\omega$), population-level cross-feeding interactions ($\sigma$ parameters), and community-level assembly outcomes. As synthetic biology advances our ability to engineer metabolic strategies with precision, theoretical frameworks like ours become essential for predicting and controlling the emergent properties of microbial ecosystems.

\section*{Materials and Methods}

\subsection*{Model formulation} We consider three microbial populations in a well-mixed chemostat environment: substrate specialists (S), metabolite specialists (M), and generalists (G). Population densities are scaled by carrying capacities, yielding dimensionless variables $s$, $m$, $g \in [0, \infty)$. The governing ordinary differential equations are:
\begin{align}
\frac{ds}{dt} &= r_S s(1 + \sigma_{SM} m + a g - s)\\
\frac{dm}{dt} &= r_M m(-1 + \sigma_{MS} s + b g - m)\\
\frac{dg}{dt} &= r_G g(d + c s + e m - g)
\end{align}
Growth rates $r_i$ have units time$^{-1}$. Net interaction parameters are defined as $a = (1-\omega)\sigma_{SG} - \omega\alpha_{SG}$, $c = (1-\omega)\sigma_{GS} - \omega\alpha_{GS}$, $b = \omega\sigma_{MG} - (1-\omega)\alpha_{MG}$, $e = \omega\sigma_{GM} - (1-\omega)\alpha_{GM}$, and $d = 2\omega - 1$, where $\sigma$ coefficients represent cooperative benefits and $\alpha$ coefficients competitive costs. The pathway parameter $\omega \in [0,1]$ controls generalist metabolic allocation.

\subsection*{Analytical methods} Equilibria were found by solving $ds/dt = dm/dt = dg/dt = 0$ algebraically. For two-species subsystems, closed-form solutions were obtained via substitution. The three-species equilibrium requires numerical root-finding. Stability was assessed via Jacobian eigenvalue analysis. The Jacobian $J$ has elements $J_{ij} = \partial(\dot{N}_i)/\partial N_j$ evaluated at equilibria. Eigenvalues were computed using standard linear algebra routines; equilibria with all $\text{Re}(\lambda) < 0$ are stable.

\subsection*{Parameter space mapping} Invasion rates were computed on $200 \times 200$ grids spanning parameter ranges. For each grid point, S-M equilibrium densities were calculated analytically, then generalist invasion fitness $\lambda_G = r_G(d + cs^*_{SM} + em^*_{SM})$ evaluated. Zero-contours of $\lambda_G$ define invasion boundaries. Contour plots and heatmaps visualize parameter landscapes.

\subsection*{Numerical integration} Time series were generated using Python's \texttt{scipy.integrate.solve\_ivp} with the Runge-Kutta RK45 method, relative tolerance $10^{-8}$, absolute tolerance $10^{-10}$. Initial conditions were drawn from biologically plausible ranges; results showed insensitivity to initial conditions given stable equilibria.

\subsection*{Parameter values} Baseline: $r_S = 1.0$, $r_M = 0.8$, $r_G = 0.9$ day$^{-1}$; $\sigma_{SM} = 0.5$, $\sigma_{MS} = 1.5$; $\sigma_{SG} = \sigma_{GS} = \sigma_{MG} = \sigma_{GM} = 0.4$; $\alpha_{SG} = \alpha_{GS} = \alpha_{MG} = \alpha_{GM} = 0.3$. These satisfy S-M stability constraints and yield coexistence windows spanning $\omega \in (0.3, 0.7)$ approximately.

\subsection*{Code availability} All analysis code, figure generation scripts, and data are available at \url{https://github.com/username/repo} and have been deposited in [repository].

\acknowledgments{We thank [names] for discussions. This work was supported by [funding agencies].}

\showacknow

\bibliography{pnas-sample}

\end{document}
